\documentclass[12pt,a4paper]{article}

% ==================== 宏包 ====================
\usepackage[UTF8]{ctex}
\usepackage{geometry}
\geometry{left=2.5cm, right=2.5cm, top=2.5cm, bottom=2.5cm}
\setlength{\headheight}{14pt}

\usepackage{booktabs}
\usepackage{longtable}
\usepackage{array}
\usepackage{enumitem}
\usepackage{listings}
\usepackage{xcolor}
\usepackage{hyperref}
\usepackage{fancyhdr}
\usepackage{tabularx}
\usepackage{makecell}
\usepackage{multicol}
\usepackage{tcolorbox}

% ==================== 样式 ====================
\hypersetup{colorlinks=true, linkcolor=black!70, urlcolor=blue!60!black}

\pagestyle{fancy}
\fancyhf{}
\fancyhead[C]{\small 数据结构课程设计\quad 第十二周周报}
\fancyfoot[C]{\thepage}
\renewcommand{\headrulewidth}{0.4pt}
\renewcommand{\footrulewidth}{0.4pt}

\definecolor{headblue}{RGB}{41, 65, 122}
\definecolor{lightgray}{RGB}{245, 245, 245}

\begin{document}

% ==================== 标题 ====================
\begin{center}
    {\LARGE\bfseries\color{headblue} 数据结构课程设计周报}\\[6pt]
    {\Large 第十二周（2026年5月19日\,--\,5月25日）}\\[12pt]
    \begin{tabular}{ll}
        \bfseries 课题名称 & 基于智能体的个性化旅游系统 \\
        \bfseries 指导教师 & 郭岗 \\
        \bfseries 小组成员 & 李佳(2024211216)、王皓轩(2024211213)、崔天睿(2024211255) \\
    \end{tabular}
\end{center}

\rule{\textwidth}{0.8pt}

% ==================== 一、本周工作概述 ====================
\section*{一、本周工作概述}

本周核心工作分为两部分：一是\textbf{真实景区数据填充}，将系统库的景区数据扩充至200条以上，同时配套补充景点节点、道路边和美食数据，使系统具备可用数据基础；二是\textbf{前后端完整联调}，将5个前端功能页面全部对接真实后端API，实现端到端数据流通。

% ==================== 二、数据增强 ====================
\section*{二、景区数据大规模填充}

\subsection*{2.1 数据来源与处理}

从公开景区名录和旅游数据集中采集了覆盖全国主要城市的景区信息，通过自动化脚本处理生成批量INSERT语句导入数据库。数据维度包括名称、类型、分类、描述、热度、评分、城市、门票价格、开放时间等信息。

数据填充后各表规模如下：

\begin{center}
\begin{tabular}{l r r}
    \toprule
    \bfseries 表名 & \bfseries 填充前 & \bfseries 填充后 \\
    \midrule
    scenic\_spots & 15 & 200+ \\
    nodes & 50 & 300+ \\
    roads & 62 & 400+ \\
    foods & 15 & 80+ \\
    users & 11 & 20+ \\
    \bottomrule
\end{tabular}
\end{center}

\subsection*{2.2 景区数据覆盖}

景区数据遵循以下分布策略：

\begin{itemize}[leftmargin=2em]
    \item \textbf{地域覆盖}：涵盖北京、上海、广州、深圳、杭州、成都、西安、武汉、重庆、南京、长沙、哈尔滨等 20+ 城市
    \item \textbf{类型覆盖}：\texttt{scenic}（自然景区、历史古迹、主题公园）+ \texttt{campus}（知名高校）
    \item \textbf{分类覆盖}：自然、人文、历史、现代、综合五类，每类占比均衡
    \item \textbf{评分分布}：评分 3.0--5.0 之间正态分布，热度 1000--50000 区间
\end{itemize}

\subsection*{2.3 地图数据同步}

为北京市内主要景区补充了详细的路网数据：

\begin{center}
\begin{tabular}{l r r l}
    \toprule
    \bfseries 区域 & \bfseries 节点数 & \bfseries 道路边数 & \bfseries 说明 \\
    \midrule
    北京邮电大学 & 15 & 25 & 新增室内节点 \\
    北京大学 & 12 & 18 & 补充未名湖周边路径 \\
    清华大学 & 12 & 18 & 水木清华至荷塘区域 \\
    颐和园 & 10 & 14 & 昆明湖周边完整路径 \\
    故宫博物院 & 12 & 16 & 东西六宫区域展开 \\
    香山公园 & 10 & 12 & 新增长路线 \\
    八达岭长城 & 8 & 10 & 新增景点 \\
    天坛公园 & 6 & 8 & 新增景点 \\
    其他景区 & 各 4--8 & 各 5--10 & 简要路网 \\
    \midrule
    \textbf{合计} & \textbf{300+} & \textbf{400+} & \\
    \bottomrule
\end{tabular}
\end{center}

每条道路包含距离（米）、拥挤度（0--1）、理想速度（米/秒）、交通工具类型四项权值，支撑 Dijkstra 多策略路径规划。

\subsection*{2.4 美食数据补充}

美食数据从 15 条扩充至 80+ 条，覆盖各景区内的热门餐厅和特色小吃：

\begin{itemize}[leftmargin=2em]
    \item 菜系涵盖京菜、川菜、粤菜、鲁菜、本帮菜、西北菜、日料、西餐、快餐、素菜、甜点、饮品等 12 类
    \item 每条美食关联所属景区和位置节点，支持 Dijkstra 最短路径排序推荐
    \item 价格区间 5--300 元，评分 3.5--5.0
\end{itemize}

% ==================== 三、前后端完整联调 ====================
\section*{三、前后端完整联调}

本周完成了所有 5 个功能页面的前后端对接，各页面联调情况如下：

\subsection*{3.1 旅游推荐（RecommendView）}

\begin{itemize}[leftmargin=2em]
    \item 对接 \texttt{/api/spots/recommend}：按热度/评分排序、分类过滤、分页
    \item 对接 \texttt{/api/spots/search}：关键词模糊搜索（Trie前缀树支撑）
    \item 景区卡片正常渲染，详情弹窗展示完整字段（描述、评分、热度、门票价格等）
\end{itemize}

\subsection*{3.2 路线规划（MapView）}

\begin{itemize}[leftmargin=2em]
    \item 对接 \texttt{/api/map/graph/:areaId}：加载对应区域的路网数据（节点坐标 + 道路边）
    \item Canvas 正确绘制节点名称和道路连线，拥挤度以色阶区分
    \item 对接 \texttt{/api/route/single}：Dijkstra 最短路径规划，路径高亮展示
    \item 对接 \texttt{/api/route/multi}：TSP 多点参观路线规划
\end{itemize}

\subsection*{3.3 场所查询（FacilityView）}

\begin{itemize}[leftmargin=2em]
    \item 对接 \texttt{/api/facilities/nearby}：按区域、节点、类别查询周边设施
    \item 返回结果按 Dijkstra 实际路径距离升序排列
\end{itemize}

\subsection*{3.4 旅游日记（DiaryView）}

\begin{itemize}[leftmargin=2em]
    \item 对接 CRUD 接口：\texttt{GET/POST/PUT/DELETE /api/diaries}
    \item 全文检索模式对接 \texttt{/api/diaries/search}（倒排索引支撑）
    \item Huffman 压缩演示对接 \texttt{/api/diaries/compress}，实时展示压缩率
\end{itemize}

\subsection*{3.5 美食推荐（FoodView）}

\begin{itemize}[leftmargin=2em]
    \item 对接 \texttt{/api/foods/recommend}：Top-K 排行榜，按菜系标签过滤
    \item 对接 \texttt{/api/foods/search}：编辑距离模糊搜索，相似度进度条展示
\end{itemize}

% ==================== 四、联调问题 ====================
\section*{四、联调过程中遇到的问题与解决方案}

\begin{center}
\begin{tabular}{p{4cm} p{4cm} p{4cm}}
    \toprule
    \bfseries 问题 & \bfseries 现象 & \bfseries 解决方案 \\
    \midrule
    大字段响应缓慢 & 200+景区列表首次加载耗时5秒以上 & 添加分页参数（page/page\_size），前端懒加载 \\
    Canvas路径超长 & 400+道路边一次性绘制卡顿 & Canvas采用虚拟化绘制，缩放只绘制可见区域 \\
    搜索无结果不友好 & 空搜索直接展示空白页 & 添加"暂无结果"空状态占位，推荐热门景区 \\
    Huffman压缩异常 & 日记正文含换行符导致JSON解析失败 & 前端发送前对换行符做转义处理 \\
    跨域问题复现 & 前端开发服务器请求后端被拦截 & 后端全局 CORS 中间件增加 OPTIONS 预检处理 \\
    \bottomrule
\end{tabular}
\end{center}

% ==================== 五、代码产出统计 ====================
\section*{五、代码产出统计}

\begin{center}
\begin{tabular}{l r}
    \toprule
    \bfseries 模块 & \bfseries 新增/修改行数 \\
    \midrule
    seed\_data.sql（数据增强脚本） & $\sim$5000 \\
    recommend\_service.h（分页/排序改造） & $\sim$100 \\
    http\_server.cpp（联调修复） & $\sim$80 \\
    App.vue / AppHeader.vue（适配改造） & $\sim$60 \\
    RecommendView.vue（联调对接） & $\sim$40 \\
    MapView.vue（Canvas优化 + API对接） & $\sim$80 \\
    FacilityView.vue（数据对接） & $\sim$30 \\
    DiaryView.vue（API对接 + 压缩演示） & $\sim$50 \\
    FoodView.vue（模糊搜索对接） & $\sim$30 \\
    \midrule
    \textbf{合计} & \textbf{$\sim$5470} \\
    \bottomrule
\end{tabular}
\end{center}

% ==================== 六、下周计划 ====================
\section*{六、下周工作计划}

\begin{enumerate}[leftmargin=2em]
    \item \textbf{课程设计报告撰写} — 按照要求完成系统设计报告初稿（含需求分析、架构设计、算法说明、测试结果）
    \item \textbf{录屏演示视频} — 录制系统功能演示视频，覆盖5个核心页面的操作流程
    \item \textbf{最终验收准备} — 整理代码、数据库、文档，确保系统可正常运行
    \item \textbf{查漏补缺} — 单元测试覆盖率提升、边缘场景处理、UI细节完善
\end{enumerate}

\end{document}
